\documentclass[11pt]{article}
\usepackage[a4paper,margin=26mm]{geometry}
\usepackage{amsmath,amssymb}
\setlength{\parindent}{0pt}
\setlength{\parskip}{7pt}
\emergencystretch=2em
\title{Most-general unifiers can have shallow images\\
\large Disproof of conjecture 00000006069}
\author{ziangni-sys}
\date{4 October 2026}
\begin{document}
\maketitle
Both versions of the source assert that the depth of a most-general
unifier is bounded below by half the variable count.
Three nontrivial equations on six distinct, actually occurring
variables disprove this lower bound. The most-general substitution
has only variable images and is idempotent after one application.

\textbf{Actual first-order terms and constraints.}
Use variables $x_0,\ldots,x_5$, a constant $c$ and one unary function
symbol $F$. Terms are recursively generated by these constructors.
A substitution $\sigma$ assigns a term to each variable and extends
to all terms by
\[
 \widehat\sigma(x_i)=\sigma(x_i),\qquad
 \widehat\sigma(c)=c,\qquad
 \widehat\sigma(F(t))=F(\widehat\sigma(t)).
\]
Consider the finite equation set
\[
 E=\{x_0=x_1,\quad x_2=x_3,\quad x_4=x_5\}.
\]
All three equations have distinct sides. Every one of the six
variables occurs in $E$, and there are no unused variables.

Define the simultaneous substitution
\[
 \theta(x_0)=x_1,\quad \theta(x_2)=x_3,\quad \theta(x_4)=x_5,
 \qquad
 \theta(x_1)=x_1,\quad\theta(x_3)=x_3,\quad\theta(x_5)=x_5.
\]
Both sides of each equation become its odd-indexed variable.
Thus $\theta$ genuinely unifies all equations.

\textbf{Universal most-generality, not just a certificate.}
Let $\tau$ be any substitution unifying $E$. Its images satisfy
\[
 \tau(x_0)=\tau(x_1),\quad
 \tau(x_2)=\tau(x_3),\quad
 \tau(x_4)=\tau(x_5).
\]
For every variable, therefore,
$\tau(x_i)=\widehat\tau(\theta(x_i))$.
Induction on terms gives the full factorization
\[
 \widehat\tau(t)=\widehat\tau(\widehat\theta(t))
 \qquad\text{for every term }t.
\]
Taking the residual substitution to be $\tau$ proves that every
unifier is an instance of $\theta$. Hence $\theta$ is a
most-general unifier in the standard substitution ordering.

\newpage
\textbf{Depth under both usual leaf conventions.}
With zero-based term height, variables and constants have height
zero and $F(t)$ has height $h(t)+1$.
All six substitution images are variables, so their maximum height
is zero. With one-based leaf height their maximum height is one.
Thus
\[
 d_0(\theta)=0,\qquad d_1(\theta)=1.
\]
The actual variable count in the constraint set is six, so the
claimed lower bound would require depth at least $6/2=3$.
Both values are strictly smaller.

Even if variable count were interpreted as the number of variables
changed by the substitution, the support is exactly
$\{x_0,x_2,x_4\}$ and has cardinality three.
The one-based depth $1$ is still less than $3/2$, and zero-based
depth is smaller still. Equivalently,
\[
 2d_1(\theta)=2<3<6.
\]
The substitution is also idempotent on all terms:
\[
 \widehat\theta(\widehat\theta(t))=\widehat\theta(t).
\]
There are no substitution chains of growing depth: every changed
variable maps directly to a fixed variable in one step.
This addresses a chain-depth interpretation as well.

\textbf{Scope.}
The counterexample targets only the explicit variable-count depth
lower bound. It does not require interpreting the undefined
conversion constant, and makes no claim about complexity of a
particular occurs-check implementation.
The equations are nontrivial and all counted variables occur in
them. Most-generality is proved for arbitrary substitutions
and arbitrary terms, not asserted from a table of image values.

\textbf{Formal correspondence.}
Lean defines an inductive first-order term type, recursive
substitution application, recursive occurring-variable sets and
term heights. The actual three-element equation set and substitution
are finite objects on \texttt{Fin 6}.
\texttt{occurring\_count} and \texttt{nontrivial\_equations}
verify complete occurrence and nontriviality.

\texttt{theta\_unifies} proves all constraints.
\texttt{factor\_variables} extracts the three equalities from an
arbitrary unifier, and \texttt{theta\_most\_general} proves the
universal factorization by term induction.
\texttt{theta\_idempotent} proves the full idempotence identity.
The final \texttt{counterexample} combines actual most-generality,
both depths, both variable counts and strict failures of the bound.

\textbf{Reproduction.}
Run \texttt{lake build} in \texttt{lean/} with Lean 4.19.0 and
public Mathlib revision
\begin{center}\small\texttt{c44e0c8ee63ca166450922a373c7409c5d26b00b}\end{center}
The project prints the relevant axiom audits.
\end{document}
